% ID: SEL-D0034-OB-E03
% ESTADO: cotejada. Segunda pasada visual y borrador compilado revisados.
% Fuente: Selectividad/matematicas_aplicadas/titular_septiembre_1_MATEMATICAS_APLICADAS.pdf
% SHA-256: c637bbde5c26c0b35b59202bd832ddc97990d3845365ce483fbf2a49832f73fe
% Página: 2; opción B; ejercicio original 3.

El aparcamiento de una sala de conciertos está completo el 85\% de los días. El 90\% de los días que el aparcamiento está completo, la sala de conciertos está llena, y el 22\% de los días que el aparcamiento no está completo, la sala de conciertos no está llena.

Elegido un día al azar,
\begin{enumerate}[label=\alph*)]
\item \textbf{(1.5 puntos)} \textquestiondown{}cuál es la probabilidad de que la sala de conciertos esté llena?
\item \textbf{(1 punto)} Si se sabe que la sala de conciertos está llena, \textquestiondown{}cuál es la probabilidad de que el aparcamiento esté completo?
\end{enumerate}
